\documentclass[11pt]{article}
\usepackage[margin=1in]{geometry}
\usepackage{booktabs}
\usepackage{caption}

\title{Workforce and Housing Pressures in the Texas Triangle:\\A Brief Regional Assessment}
\author{Policy Research Division}
\date{July 2026}

\begin{document}
\maketitle

\section{Introduction}

Texas continues to absorb a disproportionate share of national population growth,
and the pressure is most visible along the urban corridor connecting Dallas--Fort
Worth, Austin, San Antonio, and Houston. This short assessment summarizes recent
indicators of labor market tightness and housing affordability across selected
metropolitan areas. The figures reported here are drawn from the division's
internal projection model and should be treated as illustrative planning
estimates rather than official statistics.

Three themes emerge from the analysis. First, employment growth remains
concentrated in a handful of sectors, particularly professional services and
health care, while goods-producing industries have grown more slowly.
Table~\ref{tab:employment} reports sector-level employment change for the two
most recent projection years. Second, housing costs have diverged sharply across
metros: the fastest-growing regions are also those where median rents have risen
most quickly, as shown in Table~\ref{tab:housing}. Third, smaller metros outside
the urban triangle are beginning to capture spillover growth, a pattern
summarized in the unnumbered reference table at the end of
Section~\ref{sec:smaller}.

\section{Labor Market Indicators}
\label{sec:labor}

The state's labor market has remained tight through the projection window.
Unemployment in the major metros has stayed below the national average, and
employers in several sectors report persistent difficulty filling mid-skill
positions. The sectoral detail in Table~\ref{tab:employment} makes clear that
service-providing industries account for nearly all net job creation in the
model, with professional and business services leading. Goods-producing
sectors, by contrast, show modest gains driven largely by construction activity
tied to residential demand.

\begin{table}[htbp]
  \caption{Projected employment change by sector, selected Texas metros}
  \label{tab:employment}
  \centering
  \begin{tabular}{lrrr}
    \toprule
    Sector & \multicolumn{2}{c}{Jobs added (thousands)} & Share of total \\
    \cmidrule(lr){2-3}
           & Year 1 & Year 2 & (\%) \\
    \midrule
    Professional services & 61.4 & 58.9 & 27.3 \\
    Health care \& social assistance & 44.2 & 47.6 & 20.8 \\
    Leisure \& hospitality & 31.8 & 29.1 & 13.9 \\
    Construction & 23.7 & 21.4 & 10.2 \\
    Manufacturing & 12.6 & 11.3 & 5.4 \\
    All other sectors & 49.3 & 51.7 & 22.4 \\
    \bottomrule
  \end{tabular}
\end{table}

Wage growth has been broad-based but uneven. Entry-level wages in hospitality
have risen faster than the metro average, reflecting acute recruitment
difficulty, while wage gains for mid-career professional workers have moderated
as in-migration has expanded the candidate pool.

\section{Housing Affordability}
\label{sec:housing}

Housing supply has not kept pace with household formation in the largest
metros. Permitting activity slowed during the rate environment of the early
2020s and has only partially recovered. The result, visible in
Table~\ref{tab:housing}, is a widening gap between rent growth and income
growth in the fastest-growing regions. Affordability pressure is most severe
for renter households below the median income, and the model projects continued
erosion absent a substantial supply response.

\begin{table}[htbp]
  \centering
  \begin{center}
  \begin{tabular}{lrrr}
    \toprule
    Metro area & Median rent (\$) & Rent growth (\%) & Vacancy (\%) \\
    \midrule
    Frisco--Plano corridor & 1{,}847 & 9.2 & 4.1 \\
    Round Rock--Georgetown & 1{,}763 & 8.7 & 4.8 \\
    Katy--Sugar Land & 1{,}592 & 6.3 & 5.6 \\
    New Braunfels--Seguin & 1{,}438 & 7.9 & 5.2 \\
    \multicolumn{4}{l}{\emph{Reference group}} \\
    Statewide composite & 1{,}376 & 5.1 & 6.4 \\
    \bottomrule
  \end{tabular}
  \end{center}
  \caption{Rental market indicators, selected suburban submarkets}
  \label{tab:housing}
\end{table}

Local governments have responded with a mix of density reforms and fee waivers,
but the effects of these policies will not be visible in the data for several
years. In the interim, affordability strain is likely to push more households
toward smaller metros and micropolitan areas.

\section{Spillover to Smaller Metros}
\label{sec:smaller}

Regions outside the urban triangle are beginning to register measurable
spillover growth. The pattern is clearest in East Texas and the Hill Country,
where net domestic in-migration has turned positive after a decade of
near-zero flows. The reference table below lists the smaller regions flagged
by the model as high-spillover candidates, along with their projected
five-year population change.

\begin{table}[htbp]
  \centering
  \begin{tabular}{lrr}
    \toprule
    Region & Population (2025) & Projected 5-yr change (\%) \\
    \midrule
    Nacogdoches area & 8{,}214 & 4.7 \\
    Fredericksburg area & 12{,}903 & 6.1 \\
    Palestine area & 47{,}358 & 3.3 \\
    Brenham area & 36{,}541 & 5.9 \\
    \bottomrule
  \end{tabular}
  \caption*{Reference list: high-spillover candidate regions (unnumbered)}
\end{table}

These smaller regions face their own capacity constraints, particularly in
water infrastructure and broadband, and the assessment's companion volume
addresses those topics in detail.

\section{Conclusion}

The indicators assembled here point to a familiar tension: the same growth
that drives employment gains (Table~\ref{tab:employment}) is straining housing
markets (Table~\ref{tab:housing}) and beginning to redistribute population
toward smaller regions. Policymakers weighing infrastructure investments
should treat the spillover regions identified above as leading candidates for
early capacity planning.

\end{document}
